\documentclass[10pt,twocolumn]{article}

\usepackage[a4paper,margin=0.72in]{geometry}
\usepackage{amsmath,amssymb,booktabs,graphicx,microtype}
\usepackage[hidelinks]{hyperref}

\newcommand{\RZZ}{\mathrm{RZZ}}
\newcommand{\evidence}[1]{\textnormal{[#1]}}

\title{Decision-Relevant Certification of Native Fractional RZZ Gates}
\author{Zhen-Jun Liu}
\date{Development draft, 2026-11-27}

\begin{document}
\maketitle

\begin{abstract}
We study whether coherent residual error in a native fractional
$\RZZ(\theta)$ gate is harmless, harmful, or statistically inconclusive for a
fixed downstream candidate-selection decision.  The development framework
combines angle-local phase probes, finite-sample confidence sets, a
decision-induced tolerance, model diagnostics, and a three-outcome rule.  An
exploratory pilot on IBM hardware established platform feasibility but
triggered model-inadequacy diagnostics; all pilot certificates were therefore
inconclusive.  Seeded development simulations identify informative regimes
for the proposed rule while also showing weak sensitivity to small visibility
curvature.  The literature, theory, diagnostics, and power gates remain open,
so this draft makes no confirmatory hardware or novelty claim.
\end{abstract}

\section{Introduction}
Average and twirled gate metrics need not determine a downstream decision when
the residual error is coherent.  This distinction is especially relevant to
continuously parameterized native gates, for which the implemented angle is
part of the decision circuit.  Prior work supplies robust phase-estimation
ideas, branch-resolved phase probes, structured diagonal learning, and
three-outcome classification mechanics
\cite{RUDINGER2025,CHAI2026,DENG2026}.  The candidate contribution examined
here is their integration with a nonlinear decision-induced tolerance and
strict finite-sample projection for native fractional $\RZZ$ gates.  This
positioning remains preliminary because the G1 review is open.

\section{Methodology}
At each decision angle, the coherent diagonal residual is modeled as
\begin{equation}
 U_e=\exp[-i(\epsilon ZZ+aZI+bIZ)/2].
\end{equation}
Control-state differencing identifies $\epsilon$ under the stated model;
matched delays and role reversal separate the two local phases
\evidence{theory, draft}.  Four analysis phases produce a confidence rectangle
for the in-phase and quadrature contrasts.  Hoeffding bounds and a union bound
give a finite-sample circular phase set and a contrast-amplitude interval
\cite{HOEFFDING1963}.

Let $\gamma(e)$ be the signed decision gap and $m\geq0$ the harmless margin.
Given a uniformly valid interval $[L,U]$, the rule reports harmless if
$L\geq m$, harmful if $U<0$, and inconclusive otherwise.  A wrong directional
claim implies noncoverage, so its probability is at most the allocated error
level \evidence{theory, draft}.  A named harmless/unsafe pair also yields a
Bretagnolle--Huber necessary sample bound; this is pair-specific and not yet a
matching characterization of the composite optimum.

\section{Experiment Setup}
The unlocked development decision family contains two- to four-qubit
$p=1$ MaxCut instances.  The prospective two-angle study uses only the
single-edge sign-sensitive instance, with independent confidence boxes at its
two decision angles.  It evaluates 1,000 seeded trials per cell at 500, 1,000,
2,000, and 4,000 shots per phase.  Diagnostic calibration uses 2,000 trials per
scenario and a family-wise development allocation of $\alpha_{\rm diag}=0.05$.

The exploratory HW-P2 pilot used \texttt{ibm\_marrakesh} under a separate
owner-approved pilot lock.  Its evidence label is
\evidence{hardware-locked, exploratory pilot HW-P2}; it is not a G5 run.

\section{Results}
Exact development witnesses show that equal average gate fidelity can coexist
with opposite decisions on named instances \evidence{exact, development},
consistent with the known separation between average and worst-case coherent
error metrics \cite{KUENG2016}.

In the diagnostic simulation, the M0 null triggered in 0 of 2,000 trials
(95\% Wilson upper bound 0.00192).  Detection was 98.4\% for a control-flip
rate of 0.03 and 87.3\% for phase curvature 0.005, but only 22.0\% for
visibility curvature 0.010 \evidence{sim, development}.  Non-triggering is
therefore not evidence that exponential visibility decay is valid.

For the sign-sensitive prospective grid, a harmful residual $\epsilon=-0.4$
was certified in 100\% of trials at 4,000 shots per phase.  At $\epsilon=0$,
the harmless rate was 95.1\% at the same budget.  A separated harmless value
$\epsilon=0.1$ reached 100\% observed power at 2,000 shots per phase.  The true
gray-zone value $\epsilon=-0.2$ remained inconclusive at every tested budget
\evidence{sim, development}.

HW-P2 exposed native fractional $\RZZ$ and all submitted circuits passed the
stored ISA checks.  Nevertheless, the model diagnostic triggered, and all
three pilot certificates were reported as inconclusive (model-inadequate)
\evidence{hardware-locked, exploratory pilot HW-P2}.

\section{Conclusion}
The current evidence supports a reproducible development framework and a
narrow platform-feasibility statement.  It does not support a confirmatory
decision certificate.  The three-outcome rule correctly preserves
inconclusive outcomes when either the decision boundary or the model
diagnostic is unresolved.

\section{Discussion}
Directly supported findings are limited to exact calculations, named
simulations, and the exploratory HW-P2 records.  A possible cause of the HW-P2
diagnostic trigger is excited-control relaxation at large repetition counts,
but the pilot does not establish a mechanism.  Generalization is limited by
the angle-local diagonal model, selected backend snapshot, one-edge power
grid, conservative interval projection, and incomplete literature review.
Multi-edge allocation, angle-dependent alternatives, drift, diagnostic cost,
and an independent mathematical audit remain open.  G4 locking and
confirmatory hardware are prohibited until the preceding gates pass.

\bibliographystyle{abbrv}
\bibliography{../docs/literature/references}

\appendix
\section{Evidence provenance}
The authoritative claim-to-file mapping is maintained in
\texttt{docs/claim\_ledger.md}.  Generated simulation outputs record seeds and
inputs in JSON.  The manuscript is a development artifact and is not a
preregistration.

\end{document}
